\section*{Capítulo \ref{ch:b3:nuclear-physics} --- Física nuclear}

\begin{solution}{exo:b3:nuclear-physics:1}
(a) $^{12}$C: $6, 6, 12$; $^{14}$C: $6, 8, 14$; $^{235}$U:
$92, 143, 235$; $^{238}$U: $92, 146, 238$. (b) Los carbonos; los
uranios. (c) La química es la nube electrónica, que fija $Z$; la
estabilidad es el equilibrio $N$--$Z$ de dentro. (d) Isóbaros ---
conectados por la desintegración $\beta^-$: $^{14}\text{C} \to {}^{14}\text{N}
+ \text{e}^- + \bar\nu$, el tictac del reloj del radiocarbono.
\end{solution}

\begin{solution}{exo:b3:nuclear-physics:2}
(a) $\rho = Am_{\text{u}}/\tfrac43\pi r_0^3A =
3m_{\text{u}}/4\pi r_0^3 \approx \qty{2.3e17}{kg/m^3}$: $A$ se
cancela. (b) $\qty{5}{mL} \to \qty{1.2e12}{kg}$ --- mil millones de
toneladas en una cucharadita. (c) $(R/R_{\text{átomo}})^3 \approx
(4.7\times10^{-5})^3 \approx 10^{-13}$: el átomo es espacio vacío
con una mota pesada. (d) Que el volumen sea $\propto A$ significa que cada nucleón
reclama un sitio fijo y se liga solo con los vecinos que toca: la
fuerza satura en el alcance del \unit{fm}.
\end{solution}

\begin{solution}{exo:b3:nuclear-physics:3}
(a) $\Delta m = 2(1.007276 + 1.008665) - 4.001506 =
\qty{0.030376}{u}$: $B = \qty{28.3}{MeV}$. (b) $B/A =
\qty{7.07}{MeV}$. (c) $28.3\,\text{MeV}/13.6\,\text{eV} \approx
2\times10^6$: seis órdenes de magnitud --- por eso el combustible nuclear
supera al químico por un millón. (d) La $\alpha$ es un
paquete premontado y excepcionalmente barato: emitirla exporta
cuatro nucleones con un descuento de ganga de \qty{28.3}{MeV}, que ningún
nucleón suelto puede igualar.
\end{solution}

\begin{solution}{exo:b3:nuclear-physics:4}
(a) $N = N_{\text{A}}/226 = \qty{2.66e21}{}$. (b) $\lambda =
\ln2/(1600\times\qty{3.156e7}{s}) = \qty{1.37e-11}{s^{-1}}$. (c)
$\mathcal A = \lambda N = \qty{3.7e10}{Bq}$: un curio, por
construcción --- el propio gramo de los Curie. (d) Tres periodos:
$\qty{0.125}{g}$ de radio (el resto, ya radón y sus
hijos, marchando hacia el plomo).
\end{solution}

\begin{solution}{exo:b3:nuclear-physics:5}
(a) Volumen $+882$; superficie $-261$; Coulomb $-0.711\times676/3.83
= -126$; asimetría $-23.7\times16/56 = -6.8$ (en \unit{MeV}). (b)
Con el apareamiento $+12/\sqrt{56} = +1.6$: $B \approx \qty{490}{MeV}$ y
$B/A = 8.76$ frente al $8.79$ medido --- tres por mil
a partir de una gota líquida. (c) El volumen contra superficie$+$Coulomb: el
cruce de sus ritmos de crecimiento es lo que talla el máximo en el hierro.
(d) Coulomb es de largo alcance: cada protón repele a todos los demás,
$\propto Z^2$, mientras que la fuerza fuerte, que satura, solo paga
$\propto A$ --- la derrota final de los pesos pesados.
\end{solution}

\begin{solution}{exo:b3:nuclear-physics:6}
(a) $\partial/\partial Z$ de los términos de Coulomb y de asimetría
da $2a_{\text{C}}Z/A^{1/3} = 4a_{\text{A}}(A - 2Z)/A$; despéjese.
(b) $Z^* = 50.5/1.163 = 43.4$: la elección de la naturaleza en $A = 101$ es
el rutenio, $Z = 44$. (c) Sin Coulomb, $Z^* = A/2$: dos escaleras
de Fermi (\cref{ch:b3:quantum-statistics}), más baratas llenas
a igual altura. (d) El asiento en el valle del fragmento es $Z^* \approx 41$:
le faltan cinco protones --- una cascada de cinco desintegraciones $\beta^-$ lo trepa
de vuelta, y cada una convierte un neutrón y emite un electrón y un
antineutrino.
\end{solution}

\begin{solution}{exo:b3:nuclear-physics:7}
(a) $t = T_{1/2}\ln(1/0.3)/\ln2 = 5730\times1.74 \approx
\qty{e4}{}$ años: un fuego del final de la glaciación. (b) $N_{^{14}\text{C}}
= (N_{\text{A}}/12)\times10^{-12} = \qty{5.0e10}{}$; $\lambda =
\qty{3.8e-12}{s^{-1}}$: $\mathcal A \approx \qty{0.2}{Bq}$ ---
una docena de desintegraciones por minuto y por gramo, la tasa de recuento de trabajo de
todos los laboratorios de datación. (c) Tras nueve periodos, la actividad
se hunde por debajo del fondo: el reloj sigue andando, pero ya no puede
leerse. (d) Las rocas nunca respiraron el carbono atmosférico; sus
relojes son los primordiales --- uranio--plomo,
potasio--argón --- con periodos de miles de millones de años.
\end{solution}

\begin{solution}{exo:b3:nuclear-physics:8}
(a) Los \qty{8.78}{MeV} de la $\alpha$ están muy por debajo del borde de
\qty{26}{MeV}: clásicamente no puede ni salir ni haber
entrado --- y aun así el periodo es de \qty{0.3}{\micro s}. (b)
$\lambda_1 = \ln2/\qty{3e-7}{s} = \qty{2.3e6}{s^{-1}}$;
$\lambda_2 = \ln2/\qty{1.4e17}{s} = \qty{4.9e-18}{s^{-1}}$:
razón $\approx 5\times10^{23}$. (c) La tasa de efecto túnel es
$\eu^{-2G}$ con el exponente de Gamow $\propto Z/\sqrt E$:
partir $E$ por la mitad sube $2G$ en unas cincuenta y cinco unidades, y
$\eu^{55} \approx 10^{24}$ --- la escandalosa recta de
Geiger--Nuttall, salida de una sola exponencial. (d) Los protones solares también
hacen túnel: la temperatura del núcleo debe ser justo lo bastante alta para que
la cola de Maxwell por el factor de Gamow sostenga la combustión ---
el mismo acantilado, escalado desde el otro lado.
\end{solution}

\begin{solution}{exo:b3:nuclear-physics:9}
(a) La partición mueve ${\sim}235$ nucleones de \qty{7.6}{} a
${\sim}\qty{8.5}{MeV}$ de \omterm{prop:b3:nuclear-physics:binding}{energía de enlace}: $235\times0.9 \approx
\qty{200}{MeV}$. (b) $N = \qty{2.56e24}{}$ núcleos:
$E = \qty{8.2e13}{J} \approx 2000$ toneladas de petróleo --- por
kilogramo. (c) \qty{3}{GW} térmicos durante un año son
\qty{9.5e16}{J}: alrededor de \qty{1.2}{t} de $^{235}$U. (d) Los dos
fragmentos positivos salen despedidos por su propia repulsión de
Coulomb, y llevan ${\sim}\qty{170}{MeV}$ de energía cinética
que se hace calor en micrómetros --- un reactor es una tetera
cuya llama es un retroceso electrostático.
\end{solution}

\begin{solution}{exo:b3:nuclear-physics:10}
(a) $10^4$ generaciones: $(1.001)^{10^4} = \eu^{10} \approx
22000$ --- de megavatios a decenas de gigavatios en un segundo; ningún
motor mueve una barra tan deprisa. (b) Con $k - 1 < \beta$, la cadena
no puede reproducirse solo con neutrones inmediatos: cada paso de crecimiento
espera a los rezagados de ${\sim}\qty{10}{s}$, de modo que el tiempo de
generación efectivo es de una fracción de segundo o más. (c)
$(1.0005)^{10} \approx 1.005$: medio punto porcentual por segundo ---
terreno de barras y operadores. (d) Llevar $k - 1$ más allá de $\beta$
puso al reactor en el reloj inmediato de \qty{e-4}{s}: perdieron el
plazo de gracia de los neutrones retardados que hace gobernables a los reactores.
\end{solution}

\begin{solution}{exo:b3:nuclear-physics:11}
(a) $\dot m = L_\odot/c^2 = \qty{4.2e9}{kg/s}$ de masa pura;
el caudal de hidrógeno $\dot m/0.0071 \approx \qty{6e11}{kg/s}$ ---
seiscientos millones de toneladas por segundo. (b) Hidrógeno quemable
${\sim}2\times10^{30}\times0.1\times0.75 = \qty{1.5e29}{kg}$:
$t \approx \qty{2.5e17}{s} \approx 8$ mil millones de años --- el Sol
está en la mediana edad. (c) Barrera de ${\sim}\qty{1.4}{MeV}$ frente a
$k_{\text{B}}T \approx \qty{1.3}{keV}$: un déficit de un factor mil
--- solo los protones más rápidos de la cola de Maxwell, haciendo túnel
(\cref{exo:b3:nuclear-physics:8}), llegan a fusionarse; y de ahí que el Sol
arda durante gigaaños en vez de estallar. (d) La ceniza es
$^4$He --- estable, con la placidez de lo doblemente mágico; no hay
fragmentos ricos en neutrones que se desintegren por $\beta$ durante siglos.
\end{solution}

\begin{solution}{exo:b3:nuclear-physics:12}
(a) $\lambda = \ln2/\qty{21600}{s} = \qty{3.2e-5}{s^{-1}}$:
$N = \mathcal A/\lambda = \qty{1.6e13}{}$ núcleos ---
\qty{2.6}{ng}. Medicina por nanogramos. (b) Seis horas sobreviven
al escáner pero no al fin de semana: la dosis se apaga sola; una
$\gamma$ pura de \qty{140}{keV} sale del cuerpo hacia la cámara
sin el peaje de $\alpha$ o $\beta$ que quema tejidos. (c) Un electrón
y un positrón en reposo se aniquilan en \emph{dos}
fotones de \qty{511}{keV}, espalda contra espalda (energía y momento,
\cref{ch:b3:relativistic-dynamics}): cada pareja en coincidencia
define una recta que atraviesa el tumor, y miles de rectas
lo triangulan --- tomografía por leyes de conservación. (d)
\qty{2}{J/kg} calientan el tejido medio milikelvin --- térmicamente, un
sorbo de agua tibia; pero entregados como ${\sim}10^{17}$ ionizaciones
por kilogramo, cada una cortando enlaces químicos de escala \unit{eV}, resultan
letales químicamente para una célula que se divide. Los julios miden calor;
los grays miden sabotaje.
\end{solution}

\begin{solution}{pb:b3:nuclear-physics:1}
\textbf{1.} $\qty{200}{MeV} = \qty{3.2e-11}{J}$:
$20\times10^6/3.2\times10^{-11} = \qty{6.3e17}{}$ fisiones por
segundo.
\textbf{2.} $5.4\times10^{22}$ fisiones al día $\times
235/N_{\text{A}}$: unos \qty{21}{g} de $^{235}$U al día.
\textbf{3.} Carbón: $\qty{1.7e12}{J}/\qty{3e7}{J/kg} \approx 58$
toneladas al día --- una razón de masas de $\sim3\times10^{6}$, que
\emph{es} la razón de \unit{MeV} a \unit{eV} de los enlaces nucleares frente a los
químicos.
\textbf{4.} $235\times(8.5 - 7.6) \approx \qty{210}{MeV}$:
coherente, y la pequeña diferencia la contabilizan los neutrones y los
neutrinos.
\textbf{5.} Energía cinética de los fragmentos $\to$ ionización $\to$ calor,
en micrómetros y microsegundos; la conducción se lo pasa al
agua --- el reactor es una tetera que frena fragmentos.
\textbf{6.} Heredan el $N/Z \approx 1.55$ del uranio, muy
por encima del fondo del valle a su $A$: cada uno debe recorrer una cadena de
desintegraciones $\beta^-$ de vuelta abajo --- radiactividad de fábrica, por
geometría del valle.
\textbf{7.} $k$ = neutrones engendrados por neutrón, una generación
después: $0.999$ se apaga, $1.000$ se mantiene (un reactor en
funcionamiento) y $1.001$ crece.
\textbf{8.} Los neutrones lentos se demoran cerca del núcleo y el
apetito de fisión del $^{235}$U crece abruptamente a energías térmicas.
Las colisiones elásticas ceden energía más deprisa contra masas iguales ---
y el hidrógeno iguala la masa del neutrón: el agua es un moderador
hecho a medida.
\textbf{9.} El boro-10 devora neutrones ($\text{n} + {}^{10}
\text{B} \to \alpha + {}^{7}\text{Li}$): barras dentro, neutrones
comidos, $k$ abajo; barras fuera, $k$ arriba --- el acelerador.
\textbf{10.} $(1.0005)^{10^4} = \eu^{5} \approx 150$: potencias
por ciento cincuenta cada segundo --- inútil para una maquinaria.
\textbf{11.} Con el reloj retardado de \qty{0.1}{s}:
$(1.0005)^{10} \approx 1.005$, medio punto porcentual por segundo. Regla:
manténgase $|k - 1|$ muy por debajo de $\beta = 0.007$, para que los
neutrones retardados tengan siempre el voto de calidad.
\textbf{12.} Hervir se lleva el moderador: los neutrones siguen rápidos,
la fisión pasa hambre y $k$ cae --- el diseño moderado por agua
se estrangula solo (una realimentación negativa de la que carecían sus primos
de grafito).
\textbf{13.} El calor residual: la radiactividad de los fragmentos ---
la pregunta 6 de la Parte I --- obedece a periodos, no a barras de control, y
sigue rindiendo megavatios justo después de la parada.
\textbf{14.} Adelantar a la luz en el agua: $v > c/n$.
\textbf{15.} $v > 0.752c$: $\gamma = 1.51$, luego $E_{\text{k}} >
0.51\times\qty{511}{keV} \approx \qty{0.26}{MeV}$.
\textbf{16.} Los electrones $\beta$ de las desintegraciones de los fragmentos llevan
\unit{MeV}, y los $\gamma$ del núcleo patean electrones por Compton hasta
energías parecidas: ambos rebasan los \qty{0.26}{MeV} --- la piscina
está llena de electrones que cumplen el requisito.
\textbf{17.} El espectro Cherenkov se refuerza hacia las longitudes de onda
cortas (en intensidad, $\propto1/\lambda^2$): al ojo le entregan
azul y violeta --- el color es la pendiente del espectro.
\textbf{18.} El agua atenúa los $\gamma$ y los neutrones
exponencialmente; ocho metros son decenas de longitudes de semiatenuación, así que
la pasarela está a la dosis de fondo. El brillo es el agua
\emph{absorbiendo} la energía de la radiación --- prueba visible de que el
blindaje se la está comiendo.
\textbf{19.} Los fragmentos de vida corta mueren en minutos
(el atenuado); la cola de vida más larga --- el mismo calor residual de la
pregunta 13 --- mantiene un brillo tenue y una carga térmica real durante
días.
\textbf{20.} Con $T_{1/2} = \qty{66}{h}$, el padre pierde un
factor ${\sim}6$ por semana: las reservas se desintegran solas,
así que los hospitales viven de una ``vaca de tecnecio'' semanal entregada por
reactores como este.
\textbf{21.} $96/66 = 1.45$ periodos: $2^{-1.45} \approx
0.36$ --- quedan unos \qty{36}{GBq}.
\textbf{22.} Termalizados a temperatura ambiente, $\lambda =
h/\sqrt{3m_{\text{n}}k_{\text{B}}T} \approx \qty{1.5}{\angstrom}$
(\cref{exo:b3:crystalline-solids:5}): nacen casados con los
espaciados de \omterm{def:b3:crystalline-solids:lattice}{red}.
\textbf{23.} Cinco años matan todos los periodos de hasta meses:
la actividad y el calor residual caen en órdenes de magnitud, hasta que el
elemento puede viajar en un contenedor blindado en vez de en una piscina.
\textbf{24.} ${\sim}\qty{200}{kW}$ --- un centenar de calefactores
domésticos en una piscina de cien toneladas: decenas de kelvin al día en el
peor caso, retirados con holgura por convección natural --- seguridad
pasiva a fuerza de capacidad calorífica.
\textbf{25.} $B/A$ paga \qty{200}{MeV} por partición y \qty{21}{g}
al día; los neutrones retardados prestan los segundos que hacen a $k$
gobernable; el agua modera, refrigera, blinda --- y brilla
de azul precisamente porque está funcionando; y el producto que
más importa se marcha en viales para los hospitales del lunes.
\end{solution}
